\documentclass[border=10pt]{standalone}
\usepackage{fontspec}
\setmainfont{Noto Serif}[Ligatures=TeXOff]
\XeTeXgenerateactualtext=1
\newfontfamily\bengalifont{Noto Serif Bengali}[Script=Bengali, NFSSFamily=bnfam]
\usepackage[Latin,Bengali]{ucharclasses}
% group-free font switching (as in header-common.tex): robust inside braces, cells and quotes
\makeatletter
\def\BN@saved{\rmdefault}
\def\BN@bn{bnfam}
\newcommand\BNin{\edef\BN@tmp{\f@family}\ifx\BN@tmp\BN@bn\else\let\BN@saved\BN@tmp\fi\fontfamily{bnfam}\selectfont}
\newcommand\BNout{\fontfamily{\BN@saved}\selectfont}
\makeatother
\setTransitionsFor{Bengali}{\BNin}{\BNout}
% Typographic quotes and dashes are not in a ucharclasses class (class 0), so a Bangla word right
% after “ or ‘ would not trigger the Bangla font. Treat them as Latin so that every switch fires.
\XeTeXcharclass"2018=\XeTeXcharclass`A
\XeTeXcharclass"2019=\XeTeXcharclass`A
\XeTeXcharclass"201C=\XeTeXcharclass`A
\XeTeXcharclass"201D=\XeTeXcharclass`A
\XeTeXcharclass"2013=\XeTeXcharclass`A
\XeTeXcharclass"2014=\XeTeXcharclass`A
\XeTeXcharclass"2026=\XeTeXcharclass`A
\usepackage{newunicodechar}
\newunicodechar{।}{{\bengalifont\char"0964}}
\newunicodechar{॥}{{\bengalifont\char"0965}}

\usepackage[table]{xcolor}
\usepackage{booktabs,tabularx,array}
\definecolor{bapilrule}{HTML}{1F2937}
\definecolor{bapilstripe}{HTML}{F3F4F6}
\renewcommand{\arraystretch}{1.28}
\setlength{\tabcolsep}{5pt}

\begin{document}
\begin{minipage}{17cm}
{\large\bfseries সারণি ৭।}\ {\large অংশীজন নিবন্ধন}\par\vspace{7pt}
\small
\rowcolors{2}{bapilstripe}{white}
\begin{tabularx}{\linewidth}{@{}>{\hsize=1.042\hsize\raggedright\arraybackslash}X>{\hsize=1.198\hsize\raggedright\arraybackslash}X>{\hsize=0.938\hsize\raggedright\arraybackslash}X>{\hsize=1.250\hsize\raggedright\arraybackslash}X>{\hsize=0.573\hsize\raggedright\arraybackslash}X@{}}
\arrayrulecolor{bapilrule}\toprule
\rowcolor{white}\textbf{অংশীজন} & \textbf{প্রধান আগ্রহ} & \textbf{BAPIL-এর জন্য ভূমিকা} & \textbf{সম্পৃক্ততার পদ্ধতি} & \textbf{পুনরাবৃত্তি} \\
\midrule
আইসিটি বিভাগ (ডাক, টেলিযোগাযোগ ও তথ্যপ্রযুক্তি মন্ত্রণালয়) & জাতীয় এআই নীতি, ডিজিটাল রূপান্তর & সমঝোতা স্মারকের অংশীদার, স্টিয়ারিং বোর্ডের সহ-সভাপতি & যৌথ পরিকল্পনা, ত্রৈমাসিক পর্যালোচনা & মাসিক \\
মন্ত্রিপরিষদ বিভাগ & সমগ্র সরকারের সমন্বয় & মন্ত্রণালয়গুলোকে একত্রিত করা & মাইলফলকে অবহিতকরণ & ষাণ্মাসিক \\
পরিকল্পনা কমিশন & পরিকল্পনা, প্রকল্প মূল্যায়ন, পরিবীক্ষণ & P1 ও P2-এর স্বাগতিক সংস্থা; ডেটা-রক্ষক & প্রোডাক্ট ওনারশিপ, সহ-নকশা & মাসিক \\
অর্থ বিভাগ & বাজেট স্বচ্ছতা ও বিশ্লেষণ & ডেটা-রক্ষক; P2-এর ব্যবহারকারী & ডেটা চুক্তি, ব্যবহারকারী পরীক্ষা & ত্রৈমাসিক \\
লেজিসলেটিভ ও সংসদ বিষয়ক বিভাগ & প্রামাণ্য আইনি পাঠ & আইন-কর্পাসের রক্ষক & একসঙ্গে প্রবেশাধিকারের চুক্তি & ত্রৈমাসিক \\
বাংলাদেশ পরিসংখ্যান ব্যুরো & সরকারি পরিসংখ্যান ও মেটাডেটা & ডেটা-রক্ষক; P5-এর অংশীদার & মেটাডেটা সামঞ্জস্য & ত্রৈমাসিক \\
এটুআই ও বাংলাদেশ কম্পিউটার কাউন্সিল & ডিজিটাল সেবা, তথ্য বাতায়ন, ডেটা সেন্টার & প্ল্যাটফর্ম ও হোস্টিং অংশীদার & কারিগরি কার্যদল & মাসিক \\
জাতীয় উপাত্ত ব্যবস্থাপনা কর্তৃপক্ষ & ডেটা নীতি ও সম্মতি & ডেটা ব্যবহারে নিয়ন্ত্রক ও উপদেষ্টা & DPIA ও চুক্তিতে পরামর্শ & প্রয়োজন অনুযায়ী \\
বাংলাদেশ এআই হাব & দক্ষতা, স্টার্টআপ, এআই গ্রহণ & পরিপূরক কর্মসূচি & কর্মবিভাজনের সমঝোতা স্মারক, যৌথ বেঞ্চমার্ক & ত্রৈমাসিক \\
বিশ্ববিদ্যালয় ও গবেষণা প্রতিষ্ঠান & গবেষণা, শিক্ষার্থী, প্রকাশনা & স্বাগতিক প্রতিষ্ঠান, অংশীদার, টীকাকার & ফেলোশিপ, যৌথ গবেষণা & নিরবচ্ছিন্ন \\
শিল্পখাত ও স্টার্টআপ & জনস্বার্থমূলক এআই-এর বাজার & প্রযুক্তি অংশীদার, গ্রহণকারী & উন্মুক্ত আহ্বান, শিল্প দিবস & ষাণ্মাসিক \\
নাগরিক সমাজ ও গণমাধ্যম & অধিকার, স্বচ্ছতা, জবাবদিহি & ERRC সদস্য, সমালোচনামূলক পর্যালোচক & উন্মুক্ত পরামর্শ, নিবন্ধন & ত্রৈমাসিক \\
উন্নয়ন সহযোগী & ডিজিটাল জনসম্পদ, সক্ষমতা & অর্থায়নকারী ও কারিগরি উপদেষ্টা & যৌথ কর্মপরিকল্পনা, পর্যালোচনা & ষাণ্মাসিক \\
নাগরিক & সহজলভ্য তথ্য, আস্থা & চূড়ান্ত উপকারভোগী & বাংলায় জনপ্রতিবেদন, মতামতের চ্যানেল & নিরবচ্ছিন্ন \\
\bottomrule
\end{tabularx}
\par\vspace{5pt}{\footnotesize\itshape কোনো প্রতিষ্ঠান তালিকাভুক্ত করা তার সম্ভাব্য ভূমিকা বর্ণনা করে; এর মাধ্যমে কোনো চুক্তি বা সমর্থন বোঝায় না।}
\par\vspace{3pt}{\scriptsize\color{gray} উৎস: BAPIL টেকনিক্যাল ব্লুপ্রিন্ট, সংস্করণ ১.০ (২০২৬)। CC BY 4.0।}
\end{minipage}
\end{document}
